% !TEX program = xelatex
\documentclass{akizumi-note}
\akizumisetup{course=Reinforcement Learning, kicker={Reinforcement Learning\akizumisep Lecture 4}}
\title{Self-Predictive Representations}
\author{Aaron Lin}

\begin{document}
\maketitle
\blindfootnote{Notes for the ELSA Lab discussion, October 2026.}

\section{Setup}
Every method has this shape:
\[
  \vb{z}_t = f_\theta(\vb{o}_t), \qquad
  \hat{\vb{z}}_{t+1} = g_\phi(\vb{z}_t, \vb{a}_t), \qquad
  \mathcal{L}_{\text{pred}} = \E\norm[\big]{\hat{\vb{z}}_{t+1} - \sg(f_{\theta'}(\vb{o}_{t+1}))}^2 .
\]
\begin{itemize}
  \item Problem: $f_\theta \equiv \text{const}$ drives the loss to zero, i.e.\ \textbf{collapse}.
  \item Fixes: stop-grad, redundancy reduction, distribution matching.
    \begin{itemize}
      \item SIGReg pushes latents toward $\mathcal{N}(\vb{0}, \vb{I})$.
    \end{itemize}
\end{itemize}

\begin{definition}[Collapse]
An encoder $f$ \emph{collapses} if $\Var[f(\vb{o})] = 0$.
\end{definition}

\begin{restatable}[Cram\'er--Wold]{theorem}{cramerwold}\label{thm:cw}
A random vector $\vb{z}$ is $\mathcal{N}(\vb{0},\vb{I})$ iff $\vb{u}\T\vb{z} \sim \mathcal{N}(0,1)$ for every unit $\vb{u}$.
\end{restatable}
The proof is in \cref{app:proofs}.\sidenote{With \texttt{layout=margin} this note sits in the margin.}

\begin{notation}
  \notn{$\vb{o}_t$, $\vb{a}_t$}{observation and action at time $t$}
  \notn{$\vb{z}_t = f_\theta(\vb{o}_t)$}{latent state}
  \notn{$\sg(\cdot)$}{stop-gradient}
\end{notation}

\begin{algorithm}[h]
\caption{SIGReg training step}\label{alg:sigreg}
\begin{algorithmic}[1]
  \Require batch $\{\vb{o}_i\}$, directions $K$, weight $\lambda$
  \State $\vb{Z} \gets f_\theta(\vb{o})$
  \State $\mathcal{L} \gets \mathcal{L}_{\text{pred}} + \lambda\,\text{SIGReg}(\vb{Z})$
  \State $\theta \gets \theta - \eta \nabla_\theta \mathcal{L}$
\end{algorithmic}
\end{algorithm}
\Cref{alg:sigreg} uses \cref{thm:cw}.

\begin{remark}
SIGReg removes the need for an EMA target. 中文內容也可以直接混排。
\end{remark}

\begin{codeblock}[lang=Python, title=sigreg.py]
def sigreg(z, num_dirs=64):
    u = torch.randn(z.shape[1], num_dirs)   # random directions
    u = u / u.norm(dim=0, keepdim=True)
    return epps_pulley(z @ u).mean()
\end{codeblock}

\subsection{Mechanisms}
\begin{table}[h]
\centering\small
\caption{Five anti-collapse mechanisms.}
\begin{tabularx}{\linewidth}{@{}l Y@{}}
\toprule
Mechanism & Examples \\
\midrule
Stop-grad / EMA target & BYOL, SPR, MR.Q \\
Distribution matching & \textbf{SIGReg} (LeJEPA, LeWM) \\
\bottomrule
\end{tabularx}
\end{table}
See \href{https://arxiv.org}{arXiv} and Theorem~\ref{thm:x}.
\begin{lemma}\label{thm:x}
$\int_0^1 x \dd x = \tfrac12$. Code: \texttt{latexmk -pvc note.tex}.
\end{lemma}
\appendix
\section{Proofs}\label{app:proofs}
\cramerwold*
\begin{proofof}{thm:cw}
Characteristic functions agree on every line through the origin.
\end{proofof}

\listoftheorems
\end{document}
